\documentclass[11pt,a4paper]{article}
\usepackage[margin=27mm]{geometry}
\usepackage{amsmath,amssymb,amsthm,booktabs,hyperref}
\hypersetup{hidelinks}
\newtheorem{lemma}{Lemma}
\newtheorem{theorem}[lemma]{Theorem}
\newcommand{\Aut}{\operatorname{Aut}}
\newcommand{\im}{\operatorname{im}}
\title{The ternary module of the residual Steiner graph on 176 vertices}
\author{Private derivation and finite certificate: independent review pending}
\date{}
\begin{document}\maketitle

This note supplies the remaining module and full automorphism assertions
for L050, Theorem~1(c) and Proposition~4: Leemans--Rodrigues,
\emph{Linear codes with complementary duals from some strongly regular
subgraphs of the McLaughlin graph}, Math. Commun.\ 21 (2016), 239--249.
We use the source's residual Steiner construction, the shared incidence
calculation, and the separately completed exact weight enumeration.
The new finite certificate uses the same matrix-algebra and minimum-support
arguments as the companion treatment of the $120$-coordinate modules.
It does not enumerate the $3^{21}$ codewords again. Statements identifying
\emph{all} irreducible representations remain separate retained inputs.

\section{The common incidence calculation}
Let $\mathcal D$ be the fixed $S(4,7,23)$ from the shared construction,
and delete a named point $\infty$.
Let $\mathcal H$ be the $176$ heptads avoiding it. The graph $\Gamma$ on
$\mathcal H$ joins two distinct blocks whose intersection has size three.
This is the complement of the intersection-one graph in the source's
Section~4. The construction checks the Steiner axiom and the intersection
sizes, so agreement is through the actual blocks, not just the four
strongly regular parameters.

Let $H$ be the incidence matrix with the $22$ remaining points as rows and
$\mathcal H$ as columns. All code spaces below are over $F=\mathbb F_3$.
The parameters of the Steiner system give $77,21,5$ blocks through a
point, pair and triple, respectively. Consequently a remaining point lies
in $77-21=56$ residual blocks and two lie in $21-5=16$.
Thus, over the integers,
\[
 HH^{\mathsf T}=40I_{22}+16J_{22}.
\]
The adjacency matrix $A$ satisfies, over $F$,
\[
 A=H^{\mathsf T}(J_{22}-I_{22})H,\qquad
 HH^{\mathsf T}=I_{22}+J_{22}.                         \tag{1}
\]
For the first identity the diagonal entry is $49-7=42$, and off the
diagonal it is $49-|X\cap Y|$, which is zero modulo three for intersection
one and one for intersection three. This is the common incidence
factorization already used for the residual graphs on $120$ and $253$
vertices; its coefficients are displayed here to fix the precise instance.

Put $V=\{x\in F^{22}:\sum x_i=0\}$ and $C=\im A$. Since $I+J$ is
nonsingular, $H$ has full row rank. Also $J-I$ has kernel $Fj$ and image
$V$, where $j$ is the constant vector on $22$ points. Hence
\[
 C=H^{\mathsf T}V,\qquad \dim C=21.                  \tag{2}
\]
For $x,y\in V$, their image scalar product is
$x^{\mathsf T}(I+J)y=x^{\mathsf T}y$.
The space $V$ is nondegenerate because $22\ne0$ in $F$, so $C$ is LCD.
The constant vector on $176$ coordinates is in $C^\perp$: its pullback
under $H$ is $56j$, orthogonal to $V$.

For completeness, write $B=J_{176}-I_{176}-A$ for the complementary
adjacency matrix. The parameters $(176,105,68,54)$ give $A^2=2A$ over
$F$, and $AJ=0$. Therefore $AB=0$. Since $A+B=J-I$ is nonsingular
($176-1\ne0$ in $F$), their images sum to the whole space. It follows
that $\im B=\ker A=C^\perp$ and $\im A=\ker B$.
Together with (2), this proves the two complement--dual identities and
$F^{176}=C\oplus C^\perp$ asserted in the source.

\section{Two general certificate arguments}
We recall explicitly the two short arguments used for the earlier
$120$-coordinate modules. They apply to a code in any length and do not
depend on the special group name.

\begin{lemma}\label{lem:matrix}
Let $G$ act linearly on a $d$-dimensional $F$-space. If $d^2$ words in its
generator matrices are linearly independent in $M_d(F)$, the module is
absolutely irreducible.
\end{lemma}
\begin{proof}
Those words span $M_d(F)$, since that is its dimension. A $G$-invariant
subspace is stable under every word and linear combination, hence under
every endomorphism. If it contains a nonzero vector, an endomorphism can
send that vector to any prescribed vector, so the subspace is the whole
space. Linear independence of the same matrix words persists under every
field extension, proving the absolute assertion.
\end{proof}

\begin{lemma}\label{lem:support}
Let $C\leq F^n$, and suppose its complete set $\mathcal S$ of minimum
nonzero word supports is known. Form the bipartite incidence graph with
coordinate vertices and one vertex for each distinct support, giving the
two parts distinct initial colours. Repeatedly refine each colour by its
old value and the multiset of neighbour colours. If successive
individualizations of coordinates have refined cell sizes
$m_1,\ldots,m_r$, and afterward all coordinates have different colours,
then $|\Aut(C)|\leq\prod_i m_i$. Here automorphisms mean coordinate
permutations, as in L050.
\end{lemma}
\begin{proof}
Every code automorphism preserves weight and permutes $\mathcal S$, so it
acts on this coloured incidence graph. Every refinement is preserved by
automorphisms fixing the previously individualized vertices. The orbit of
the next chosen coordinate is therefore contained in its current colour
cell. The orbit--stabilizer formula bounds each successive stabilizer
index by $m_i$. At the end all coordinates are fixed. The support vertices
have distinct coordinate neighbourhoods, so they too are fixed. Multiplying
the indices gives the stated upper bound. No claim that a colour cell is
an orbit, or that every support automorphism preserves the code, is needed.
\end{proof}

\section{The exact module and group certificate}
The literal permutations in \texttt{module\_data.json} preserve every
adjacency entry of $\Gamma$. Their closure, computed by exact permutation
composition, has order $443520$. They also permute the $22$ point-incidence
rows of $H$. The resulting faithful degree-$22$ action is explicitly
identified with \texttt{MathieuGroup(22)} by an isomorphism checked in GAP.
Thus the subgroup $G$ they generate is $M_{22}$. This conventional group
identification is distinct from proving that $G$ is the \emph{full} code
automorphism group. The latter will follow from Lemma~\ref{lem:support}.
The auxiliary group calculation also verifies that the degree-$176$ action
is primitive with stabilizer $A_7$ of order $2520$.

The row-reduced adjacency basis is checked byte for byte against the basis
of the retained complete enumeration. For each coordinate generator, its
images on these $21$ rows give a $21\times21$ representation matrix.
The checker reconstructs those images from the proposed matrices, so the
matrix module is the actual code $C$, not a representation with only the
same dimension. Deterministic multiplication and elimination produce
$441$ independent generator words, all of length at most eight. Each word
is evaluated again and the $441$ flattened matrices are row-reduced.
Their rank is $441$. A separate GAP calculation confirms the same generated
algebra dimension. Lemma~\ref{lem:matrix} proves that $C$ is absolutely
irreducible. This action is faithful as well: the $176$ columns of a code
basis are distinct, so a coordinate permutation fixing every codeword
fixes every coordinate. Column distinctness is checked in the certificate.

The retained complete enumeration gives $d(C)=56$ and exactly $924$ words
of weight $56$. The supplied literal weight-$56$ word belongs to $C$.
Its orbit under $G$ has $462$ words and is disjoint from the orbit of its
negative. Together those two orbits have $924$ words, so they exhaust the
minimum words. Their supports are the same $462$ distinct subsets.
These assertions are checked by finite orbits and the already accepted
count; no new exhaustive codeword computation is involved.

For this support-incidence graph, with $176+462$ vertices, the deterministic
refinement rule of Lemma~\ref{lem:support} successively individualizes
coordinates $1,5,29,61,2$. The corresponding cell sizes are
\[
 176,\quad70,\quad6,\quad2,\quad3.
\]
All coordinate colours then differ. Therefore
\[
 |\Aut(C)|\leq176\cdot70\cdot6\cdot2\cdot3=443520.
\]
Every graph automorphism preserves its adjacency row space, so
$G\leq\Aut(\Gamma)\leq\Aut(C)$. The already checked order of $G$ forces
equality throughout:
\[
 \Aut(C)=\Aut(\Gamma)=G\cong M_{22}.                 \tag{3}
\]
This upper bound uses no graph-isomorphism backend. Such software was
used only to discover the two literal generators.

The complete retained weight distribution also determines the dual one
by the classical MacWilliams identity
\[
 A_j^\perp=3^{-21}\sum_{i=0}^{176}A_i
 \sum_s(-1)^s2^{j-s}\binom{i}{s}\binom{176-i}{j-s}.
\]
The certificate computes all coefficients exactly, checks nonnegative
integrality and total $3^{155}$, and obtains zero at weights one through
five and $184800$ at weight six. Thus the parameters are
$[176,21,56]_3$ and $[176,155,6]_3$, with the asserted minimum counts.

\section{The precise source boundary}
Theorem~1(c) describes the \emph{unique $21$-dimensional} irreducible
$F_3$-module for $M_{22}$. Proposition~4's final sentence omits the dimension,
but its proof immediately restores it and invokes the known $3$-modular
character table. We retain that intended classification statement, and
the proof's further assertion that $21$ is the smallest nontrivial
irreducible dimension, from the source's reference [13]:
Jansen--Lux--Parker--Wilson, \emph{An Atlas of Brauer Characters} (1995).
These assertions are understood up to module isomorphism. Our full
matrix-algebra calculation proves irreducibility of this particular
module; it does not classify all modules. The unqualified literal
uniqueness sentence cannot be correct, since the trivial one-dimensional
module is irreducible too.

The source's historical near-optimality comparison for the dual is a
separate assertion based on its tables, not a current table claim or a
consequence of (3). No additional irreducibility is asserted for the
$155$-dimensional dual, which contains the constant line. The group in
(3) consists of coordinate permutations; allowing scalar monomial maps
would change the automorphism convention.

The complete mandatory local files are the shared Steiner constructor,
the literal data, the certificate and the GAP identification, together
with the retained exact enumeration data and its separately reviewed
method. The orchestration driver and its weakening controls are listed
separately. No Lean formalization or whole-L050 admission is claimed in
this bounded supplement, and no software-inclusive compression figure
is inferred from it.
\end{document}
